% =====================================================================
% 6. SONUÇ (Conclusion)
% Summarize actual findings only. Introduce NO new results.
% Evidence base: docs/FINAL_MODEL.md, docs/results_progress.md, docs/decisions.md.
% =====================================================================
\section{Sonuç}
\label{sec:conclusion}

Bu çalışmada, HiperKvasir 23-sınıf gastrointestinal endoskopi görevi için üç
önceden eğitilmiş CNN omurga ağının (ResNet50, MobileNetV2, EfficientNetB0)
özelliklerini birleştiren, MLP sınıflandırıcılı bir füzyon mimarisi, resmi
5-katlı protokol altında kontrollü biçimde değerlendirilmiştir. Karşılaştırmalar
omurga ağ, füzyon (birleştirme / ağırlıklı / GMU) ve transfer öğrenme (frozen /
fine-tune) eksenlerinde yürütülmüş; ana ölçüt olarak makro-F1 kullanılmıştır.
Resmi 5-katlı protokolde en iyi yapılandırma üçlü ağırlıklı fine-tune (exp 11)
olmuş; nihai model bu yapılandırmanın TTA'lı sürümüdür (havuzlanmış
makro-F1 $0{,}6075$). Sayısal sonuçlar ve yorumlar
Bölüm~\ref{sec:results}--\ref{sec:discussion}'de verilmiştir.

\paragraph{Sınırlamalar.}
En iyi yapılandırmanın üzerinde değerlendirilen toplamsal ablasyonların --- focal
loss, TTA ve sızıntısız seed topluluğu --- hiçbiri \%95 güven aralığı düzeyinde
en iyi çapraz-entropi yapılandırmasını istatistiksel olarak geçememiştir; bu, bu
ek tekniklerle elde edilebilecek iyileşmenin sınırlı kaldığını gösterir. Kalan açığın başlıca kaynağı, çok düşük destekli
(support $\le 10$) nadir sınıflardır; bu, kayıp/çıkarım/topluluk düzeyindeki
ayarlamalarla kapanmayan bir veri kıtlığı sınırıdır. Ayrıca, farklı protokol ve
ön işleme kullanan literatür sonuçlarıyla doğrudan baş-başa karşılaştırma
yapılamaz; yalnızca resmi-bölme tabanlı sonuçlar bağlamsal bir çapa sağlar. Son
olarak, final model tek tohumla (seed 42) eğitilmiştir.

\paragraph{Gelecek çalışma.}
İleride, nihai model üzerinde yorumlanabilirlik analizleri --- Grad-CAM++
\cite{chattopadhyay2018gradcampp} ile sınıf-etkin bölge görselleştirmesi ve
birleşik 512-boyutlu özelliklerin UMAP~\cite{mcinnes2018umap} ile izdüşürülmesi
--- eklenebilir. Eğitim reçetesinin daha ileri bir sürümü (recipe-v2) ve daha
geniş bir tohum topluluğu da, özellikle nadir sınıflar için, olası iyileştirme
yönleri olarak değerlendirilebilir. Bu yönler bu raporun kapsamı dışındadır ve
gelecek çalışma olarak bırakılmıştır.
